% Abhakseite „Das kann ich“ – Zone nur im Gesamt (\mitzone)
%% TODO Ich-kann-Sätze: Platzhalter wie in den Aufgabendateien
\begin{abhakseite}
\ifdefined\mitzone
\abhakgruppe{Kennst du schon}
\abhak{1}{Kumulierte Binomialwahrscheinlichkeiten berechnen und lesen}
\abhak{2}{Gegenereignis und Komplementärsummen („mindestens“ gegen „höchstens“)}
\abhak{3}{Erwartungswert n · p als Orientierung, auf welcher Seite der Ablehnungsbereich liegt}
\abhak{4}{Funktionsgraphen ablesen und deuten (Werte, Monotonie)}
\abhak{5}{Den Rechner für kumulierte Verteilungswahrscheinlichkeiten einsetzen}
\abhak{6}{Gegenereignis und Komplementärsummen („mindestens“ gegen „höchstens“) – Fehler finden}
\abhak{7}{Gegenereignis und Komplementärsummen („mindestens“ gegen „höchstens“) – selbst rechnen}
\fi
\abhakgruppe{Einheit 1 · Entscheidungsregel bestimmen: die Nullhypothese liefert p und die Binomialverteilung der Testgröße, die Alternative die Seite des Ablehnungsbereichs (rechtsseitig bei „mehr als“, linksseitig bei „weniger als“); die Grenze über kumulierte Wahrscheinlichkeiten so wählen, dass das Signifikanzniveau eingehalten wird}
\abhak{8}{Entscheidungsregel}
\abhak{9}{Ablehnungsgrenze bei größerem Stichprobenumfang ohne höheren Fehler erster Art bestimmen}
\abhak{10}{Entscheidungsregel – Fehler finden, Begründen, Anwendung}
\abhakgruppe{Einheit 2 · Die Nullhypothese wählen: aus der Sicht des Entscheiders}
\abhak{11}{Wahl der Nullhypothese}
\abhak{12}{Wahl der Nullhypothese – Fehler finden, Begründen, Anwendung}
\abhakgruppe{Einheit 3 · Fehlerarten und Güte: den Fehler zweiter Art für selbst gewählte Anteile berechnen (p dort wählen, wo die Nullhypothese falsch ist), einordnen und im Sachzusammenhang beschreiben; Mindestanteile für eine Fehlerschranke; die Gütekurve lesen (Ablehnwahrscheinlichkeit in Abhängigkeit von p}
\abhak{13}{Fehlerarten und Güte}
\abhak{14}{Fehlentscheidungen eines Tests im Sachzusammenhang beschreiben}
\abhak{15}{Fehlerarten und Güte – Fehler finden, Begründen, Anwendung, Darstellungswechsel}
\end{abhakseite}
